\chapter{Conclusion and Future Work}

\section{Conclusion}

The Institutional Recruitment Portal has been developed to digitize and streamline the recruitment process for academic projects. By replacing manual, email-based workflows with a secure MERN-stack platform, the system significantly improves organizational efficiency for both professors and administrative staff.

The system features a robust, role-based architecture that caters to Applicants, Professors, and Administrators. Key technical achievements include the implementation of a structured local file storage system, bulk data export capabilities (ZIP and CSV), and a secure authentication framework. The use of Docker ensures that the system is easily deployable and scalable within an institutional IT environment.

Evaluation of the portal shows that it effectively addresses the primary challenges of application fragmentation and manual document management. It provides a professional and transparent interface for applicants while giving professors the tools they need to manage recruitment at scale.

\section{Limitations}

Despite its functionality, the current version has some limitations:

\begin{enumerate}
  \item \textbf{Manual Approval for Professors:} Currently, professors must be manually approved by an admin, which can cause delays during peak recruitment periods.
  
  \item \textbf{Basic Search Functionality:} The project search is based on simple keyword matching and could be improved with more advanced filtering.
  
  \item \textbf{Lack of Real-time Notifications:} The system does not currently notify users of status changes via email or SMS in real-time.
  
  \item \textbf{No Automated Screening:} Professors still need to manually review every application; there is no automated ranking or screening based on qualifications.
\end{enumerate}

\section{Future Work}

Several areas have been identified for future enhancement:

\begin{itemize}
  \item \textbf{AI-Based Applicant Ranking:} Integrating a machine learning module to rank applicants based on their skills and credentials relative to project requirements.
  
  \item \textbf{Email and SMS Integration:} Adding automated email notifications for application submission confirmations and status updates.
  
  \item \textbf{Integrated Interview Module:} Developing a built-in scheduling and video conferencing tool for conducting online interviews within the portal.
  
  \item \textbf{Advanced Analytics:} Implementing a more detailed analytics dashboard for admins to track recruitment trends and diversity metrics.
  
  \item \textbf{Public API:} Developing a documented REST API to allow integration with other institutional ERP systems.
\end{itemize}

In conclusion, the Institutional Recruitment Portal establishes a solid digital foundation for project-based hiring. It is a production-ready solution that can be further expanded to become a comprehensive human resource management tool for academic institutions.